\slcol{ಅರ್ಜುನ ಉವಾಚ ।\\
\Index{ಕಿಂ ತದ್ಬ್ರಹ್ಮ ಕಿಮಧ್ಯಾತ್ಮಂ} ಕಿಂ ಕರ್ಮ ಪುರುಷೋತ್ತಮ ।\\
ಅಧಿಭೂತಂ ಚ ಕಿಂ ಪ್ರೋಕ್ತಮಧಿದೈವಂ ಕಿಮುಚ್ಯತೇ ॥ ೧ ॥}
\cquote{ಅರ್ಜುನನು ಹೇಳಿದನು,
 ಪುರುಷೋತ್ತಮ, ಆ ಬ್ರಹ್ಮವೆಂದರೇನು? ಆಧ್ಯಾತ್ಮವೆಂದರೆ ಏನು? ಕರ್ಮವೆಂಬುದೇನು? ಅಧಿಭೂತ ಎಂಬುದು ಯಾವುದು? ಅಧಿದೈವ ಎಂದು ಯಾವುದಕ್ಕೆ ಹೇಳುತ್ತಾರೆ?}
\slcol{\Index{ಅಧಿಯಙ್ಞಃ ಕಥಂ} ಕೋಽತ್ರ ದೇಹೇಽಸ್ಮಿನ್ಮಧುಸೂದನ ।\\
ಪ್ರಯಾಣಕಾಲೇ ಚ ಕಥಂ ಜ್ಞೇಯೋಽಸಿ ನಿಯತಾತ್ಮಭಿಃ ॥ ೨ ॥}
\cquote{ಆದಿ ಯಜ್ಞವೆಂದರೆ ಈ ದೇಹದಲ್ಲಿ ಯಾರು? ಮನಸ್ಸನ್ನು ಬಿಗಿಹಿಡಿದವರು ಸಾವಿನ ಸಮಯದಲ್ಲಿ ಕೂಡ ಹೇಗೆ ನಿನ್ನನ್ನು ತಿಳಿಯಲಾಗುವುದು?}


\newpage
\begin{mananam}{\mananamTitle{ಮನನ ಶ್ಲೋಕ - ೧, ೨}}
\nfntmtxt ಈ ಬ್ರಹ್ಮಾಂಡದ ಸ್ವರೂಪ, ಮೂಲ ಮತ್ತು ಅದನ್ನು ಸುಸ್ಥಿರವಾಗಿಡುವ ಶಕ್ತಿಗಳ ಬಗ್ಗೆ ನಾನು ಎಂದಾದರೂ ಆಳವಾಗಿ ವಿಚಾರ ಮಾಡಿದ್ದೇನೆಯೇ  ಅಥವಾ, ನನಗೆ ಹಸ್ತಾಂತರಿಸಿದ  ಉತ್ತರಗಳಿಂದ  ನನ್ನ ಕುತೂಹಲಗಳು ತಣಿಸಲ್ಪಟ್ಟು, ಈಗ ಅಂತಹ ವಿಚಾರಗಳ ಬಗ್ಗೆ ಯೋಚಿಸುವುದು ಅಪ್ರಾಯೋಗಿಕ  ಮತ್ತು ಸಮಯ ವ್ಯರ್ಥ ಎಂದು  ಪರಿಗಣಿಸಲು ಕಾರಣವಾಯಿತೇ ? 
ಜೀವನದ ಸ್ವರೂಪದ ಕುರಿತು ನನ್ನ ವಿಚಾರಣೆಗಳೇನು? ಜೀವನದಲ್ಲಿ ಇನ್ನಷ್ಟು ಆಳವಾದ ಅರ್ಥ ಮತ್ತು ಉದ್ದೇಶವನ್ನು ಕಾಣಲು ನಾನು ಇನ್ನೂ ಆಶಿಸುತ್ತೇನೆಯೇ? ಮರಣಾನಂತರದ ಜೀವನ ಮತ್ತು ಪುನರ್ಜನ್ಮದ ಬಗ್ಗೆ ನನ್ನ ಅನಿಸಿಕೆಗಳು ಮತ್ತು ನಂಬಿಕೆಗಳು ಯಾವುವು?
\end{mananam}
\sReflect
\begin{inspiration}{\mananamTitle{ಸ್ಪೂರ್ತಿ}}
\nfntmtxt ವಿಜ್ಞಾನದಲ್ಲಾಗಲಿ  ಅಥವಾ ಆಧ್ಯಾತ್ಮಿಕ ಕ್ಷೇತ್ರದಲ್ಲಾಗಲಿ, ಸರಿಯಾದ ಪ್ರಶ್ನೆಗಳನ್ನು ಕೇಳುವುದರಿಂದ ಎಲ್ಲಾ ಆವಿಷ್ಕಾರಗಳು ಪ್ರಾರಂಭವಾಗುತ್ತವೆ. ವಿಚಾರಣೆಗಳಲ್ಲಿ ತೊಡಗುವುದು ಮತ್ತು ಆಳವಾಗಿ ಪರಿಶೋಧಿಸುವುದು ಕೂಡ, ತತ್ವಜ್ಞಾನಿಗಳು ಹಾಗೂ ಋಷಿಮುನಿಗಳು ಅಭ್ಯಾಸ ಮಾಡಿದ ಒಂದು ಚಿಂತನೆಯ ಕಲೆಯೇ ಆಗಿದೆ. ಸತ್ಯಕ್ಕಾಗಿ ಬುದ್ಧನ ಅನ್ವೇಷಣೆಯು ಪ್ರಾರಂಭವಾದುದು ಸಾಮಾನ್ಯ ಸ್ವರೂಪದ ದುಃಖದ ಬಗ್ಗೆ ವಿಚಾರಣೆಯಿಂದ.  “ನಾನು ಯಾರು?” ಎಂಬ ಚಿಂತನೆಯಿಂದ ರಮಣಮರ್ಷಿಗಳಿಗೆ ಸಾಕ್ಷಾತ್ಕಾರವಾಯಿತು. 
ಧರ್ಮಗ್ರಂಥಗಳನ್ನು ಅಧ್ಯಯನ ಮಾಡುವ ಮೌಲ್ಯವು ಸಿದ್ಧ ಉತ್ತರಗಳ ಸಂಗ್ರಹವನ್ನು ಅರಸುವುದಲ್ಲ ಆದರೆ, ವೈಯಕ್ತಿಕ ಸಾಕ್ಷಾತ್ಕಾರಕ್ಕೆ ಕಾರಣವಾಗುವ ಸರಿಯಾದ ಪ್ರಶ್ನೆಗಳನ್ನು ಹೇಗೆ ಕೇಳಬೇಕೆಂದು ಕಲಿಯುವುದು. ಇದೇ ಜೀವನದಲ್ಲಿಯೇ ಉತ್ತರಗಳನ್ನು ಕಂಡುಕೊಳ್ಳಬೇಕೆಂಬ ತುರ್ತುಭಾವದೊಂದಿಗೆ, ವಿಚಾರಣೆಯ ಜ್ಯೋತಿಯನ್ನು ನಾವು ನಿರಂತರವಾಗಿ ಪ್ರಜ್ವಲಿಸುತ್ತಿರಲೇಬೇಕು!
\end{inspiration}
\newpage

\slcol{ಶ್ರೀಭಗವಾನುವಾಚ ।\\
\Index{ಅಕ್ಷರಂ ಬ್ರಹ್ಮ ಪರಮಂ} ಸ್ವಭಾವೋಽಧ್ಯಾತ್ಮಮುಚ್ಯತೇ ।\\
ಭೂತಭಾವೋದ್ಭವಕರೋ ವಿಸರ್ಗಃ ಕರ್ಮಸಂಙ್ಞಿತಃ ॥ ೩ ॥}
\cquote{ಭಗವಂತನು ಹೇಳಿದನು, ನಾಶರಹಿತವಾದ ಪರಬ್ರಹ್ಮವೇ ಇಲ್ಲಿ ಹೇಳಿದ ಬ್ರಹ್ಮ. ಪ್ರತಿ ದೇಹದಲ್ಲಿಯೂ ಆತ್ಮವಾಗಿರುವುದು ಆಧ್ಯಾತ್ಮವೆನಿಸುವುದು. ಜೀವ ಜಡಗಳ ಸೃಷ್ಟಿಗೆ ಕಾರಣವಾದ ಭಗವಂತನ ಕ್ರಿಯೆಯೇ ನಿಜವಾದ ಕರ್ಮವೆನಿಸುವುದು.}
\slcol{\Index{ಅಧಿಭೂತಂ ಕ್ಷರೋ ಭಾವಃ} ಪುರುಷಶ್ಚಾಧಿದೈವತಮ್ ।\\
ಅಧಿಯಙ್ಞೋಽಹಮೇವಾತ್ರ ದೇಹೇ ದೇಹಭೃತಾಂ ವರ ॥ ೪ ॥}
\cquote{ಅರ್ಜುನ, ನಾಶವಾಗುವ ವಸ್ತುಗಳೆಲ್ಲ ಅಧಿಭೂತ. ವಿರಾಟ್ ಪುರುಷನಾದ ಹಿರಣ್ಯ ಗರ್ಭನು ಅಧಿದೈವತ. ಈ ಶರೀರದಲ್ಲಿ ಅಧಿಯಜ್ಞನು ನಾನೇ.}
\slcol{\Index{ಅಂತಕಾಲೇ ಚ ಮಾಮೇವ} ಸ್ಮರನ್ಮುಕ್ತ್ವಾ ಕಲೇವರಮ್ ।\\
ಯಃ ಪ್ರಯಾತಿ ಸ ಮದ್ಭಾವಂ ಯಾತಿ ನಾಸ್ತ್ಯತ್ರ ಸಂಶಯಃ ॥ ೫ ॥}
\cquote{ಯಾವನು ಸಾಯುವಾಗಲು ಕೂಡ ನನ್ನನ್ನೇ ಧ್ಯಾನಿಸುತ್ತಾ ಶರೀರವನ್ನು ಬಿಟ್ಟು ಹೋಗುತ್ತಾನೋ ಅವನು ನನ್ನನ್ನು ಪಡೆಯುತ್ತಾನೆ. ಈ ವಿಷಯದಲ್ಲಿ ಸಂದೇಹವಿಲ್ಲ.}
\slcol{\Index{ಯಂ ಯಂ ವಾಪಿ} ಸ್ಮರನ್ಭಾವಂ ತ್ಯಜತ್ಯಂತೇ ಕಲೇವರಮ್ ।\\
ತಂ ತಮೇವೈತಿ ಕೌಂತೇಯ ಸದಾ ತದ್ಭಾವಭಾವಿತಃ ॥ ೬ ॥}
\cquote{ಕುಂತೀಪುತ್ರ, ಎಡಬಿಡದೆ ಚಿಂತಿಸಿದ ಸಂಸ್ಕಾರದ ಬಲದಿಂದ ಕೊನೆಯ ಕಾಲದಲ್ಲಿ ಯಾವ ಯಾವ ರೂಪವನ್ನು ಚಿಂತಿಸುತ್ತಾ ಶರೀರವನ್ನು ಬಿಡುವನೋ ಆಯಾ ಗತಿಯನ್ನೇ ಪಡೆಯುತ್ತಾನೆ.}
\slcol{\Index{ತಸ್ಮಾತ್ಸರ್ವೇಷು ಕಾಲೇಷು} ಮಾಮನುಸ್ಮರ ಯುಧ್ಯ ಚ ।\\
ಮಯ್ಯರ್ಪಿತಮನೋಬುದ್ಧಿರ್ಮಾಮೇವೈಷ್ಯಸ್ಯಸಂಶಯಮ್ ॥ ೭ ॥}
\cquote{ಆದ್ದರಿಂದ ಎಲ್ಲಾ ಸಮಯಗಳಲ್ಲಿಯೂ ನನ್ನನ್ನು ಚಿಂತಿಸು, ಯುದ್ಧವನ್ನು ಮಾಡು, ಮನಸ್ಸನ್ನು ಬುದ್ಧಿಯನ್ನು ನನಗೆ ಒಪ್ಪಿಸಿದವನಾಗಿ ನನ್ನನ್ನೇ ಹೊಂದುತ್ತಿ, ಈ ಮಾತಿಗೆರಡಿಲ್ಲ.}

\newpage
\begin{mananam}{\mananamTitle{ಮನನ ಶ್ಲೋಕ - ೫, ೬}}
\nfntmtxt ಸಾವಿನ ಕ್ಷಣದಲ್ಲಿ ಮನಸ್ಸಿನ ಸ್ಥಿತಿಯ ಮಹತ್ವವನ್ನು ನಾನು ಎಂದಾದರೂ ಪರಿಗಣಿಸಿದ್ದೇನೆಯೇ?  ಅಥವಾ, ನಾನು ಕೇವಲ, ಜೀವನವನ್ನು ಆನಂದಿಸುವತ್ತ ಮತ್ತು ಪ್ರತಿದಿನ ಒಳ್ಳೆಯ ಅನುಭವಗಳನ್ನು ಹೊಂದುವತ್ತ ಮಾತ್ರ ಗಮನಹರಿಸಿದ್ದೇನೆಯೇ? ನಿದ್ರೆಗೆ ಮುಂಚಿತವಾಗಿ ನನ್ನ ಆಲೋಚನೆಗಳ ಸ್ವರೂಪ ಮತ್ತು ಅವು ನನ್ನ ಕನಸುಗಳ ಮೇಲೆ ಹೇಗೆ ಪ್ರಭಾವ ಬೀರುತ್ತವೆ ಎಂಬುದರ ನಡುವಿನ ಸಂಪರ್ಕವನ್ನು ನಾನು ಅನುಭವಿಸಿದ್ದೇನೆಯೇ? 
ನನ್ನ ದೇಹವನ್ನು ತೊರೆಯುವ ಕ್ಷಣದಲ್ಲಿ ನನ್ನ ಆಲೋಚನೆಗಳು,  ದೇವರ ಮೇಲೆ ಅಥವಾ ನನ್ನೊಳಗಿರುವ ದೈವತ್ವದ ಮೇಲೆಯೇ ಕೇಂದ್ರೀಕರಿಸಿದೆ ಎಂದು ನಾನು ಹೇಗೆ ಖಚಿತಪಡಿಸಿಕೊಳ್ಳಬಹುದು? ಜೀವನದ ಕಠಿಣ ಕ್ಷಣಗಳಲ್ಲಿ, ನನ್ನ ಆಲೋಚನೆಗಳು ಭಗವಂತನತ್ತ ತಿರುಗುತ್ತವೆಯೇ?
\end{mananam}
\sReflect
\begin{inspiration}{\mananamTitle{ಸ್ಪೂರ್ತಿ}}
\nfntmtxt \footnotesize ನಮ್ಮ ಎಲ್ಲಾ ಆಲೋಚನೆಗಳು ಮತ್ತು ಪ್ರಯತ್ನಗಳು ಮನಸ್ಸಿನ ಅಗೋಚರ ಅಂತರಂಗದಲ್ಲಿ ಗುರುತುಗಳನ್ನು ಮೂಡಿಸುತ್ತವೆ; ಅವೇ ಸಂಸ್ಕಾರಗಳಾಗಿ ನಮ್ಮ ಎಲ್ಲಾ ಚಟುವಟಿಕೆಗಳಿಗೆ ಕಾರಣವಾಗುತ್ತವೆ.
ನಮ್ಮ ಮುಂದಿನ ಜೀವನ (ಯಾವುದೇ ಲೋಕದಲ್ಲಿ ಪುನರ್ಜನ್ಮವಾಗಿರಬಹುದು) ಅತ್ಯಂತ ಪ್ರಬಲ ಹಾಗೂ ಪುನಃ ಪುನಃ ದೃಢೀಕರಿಸಲ್ಪಟ್ಟ ಸಂಸ್ಕಾರಗಳ ಪ್ರಭಾವಗಳಿಂದ ನಿರ್ಧರಿಸಲ್ಪಡುತ್ತದೆ ಎಂದು ಧರ್ಮಗ್ರಂಥಗಳು ಪ್ರತಿಪಾದಿಸುತ್ತವೆ.  
ವಾಸ್ತವವಾಗಿ ಇನ್ನೂ ಕೆಲ ಧರ್ಮ ಗ್ರಂಥಗಳು, ಈ ದೇಹವನ್ನು ಬಿಟ್ಟು ಹೋಗುವ ನಿರ್ಣಾಯಕ ಕ್ಷಣದಲ್ಲಿ ‘ಕೇವಲ ಪುನರ್ಜನ್ಮವಲ್ಲ ಆದರೆ, ವಿಮೋಚನೆಯನ್ನೂ ಸಾಧಿಸಬಹುದು’ ಎಂದು ಸೂಚಿಸುತ್ತವೆ. ಸಾವಿನ ಸಮಯದಲ್ಲಿ ನಮ್ಮ ಎಲ್ಲಾ ಆಲೋಚನೆಗಳೂ ಆ ಒಂದು ದೈವಿಕ ತತ್ವವನ್ನೇ ಹಿಡಿದಿಟ್ಟುಕೊಳ್ಳಬಹುದಾದರೆ  ಅದು ಮೋಕ್ಷಕ್ಕೆ ಅತ್ಯಂತ ಪರಿಣಾಮಕಾರಿಯಾದ ಮಾರ್ಗವಾಗುತ್ತದೆ. 
ಆದಾಗ್ಯೂ, ಒಬ್ಬರು ತಮ್ಮ ಜೀವನದುದ್ದಕ್ಕೂ ಅದನ್ನು ಅಭ್ಯಾಸ ಮಾಡಿದ್ದಲ್ಲಿ ಮಾತ್ರ, ಸಾವಿನ ಕ್ಷಣದಲ್ಲಿ ಯಾವುದೇ ಅಪೇಕ್ಷಿತ ವಿಷಯದ ಬಗ್ಗೆ  ಆಲೋಚಿಸಬಹುದು. ಹೀಗಾಗಿ, ಒಬ್ಬ ಸಾಧಕನಿಗೆ ದೈವಿಕ ಉಪಸ್ಥಿತಿಯಲ್ಲಿರುವುದನ್ನು ಅಭ್ಯಾಸ ಮಾಡಲು ಇಡೀ ಜೀವನವೇ ಒಂದು ಅವಕಾಶವಾಗಿರುತ್ತದೆ. 
\end{inspiration}
\newpage

\slcol{\Index{ಅಭ್ಯಾಸಯೋಗಯುಕ್ತೇನ} ಚೇತಸಾ ನಾನ್ಯಗಾಮಿನಾ ।\\
ಪರಮಂ ಪುರುಷಂ ದಿವ್ಯಂ ಯಾತಿ ಪಾರ್ಥಾನುಚಿಂತಯನ್ ॥ ೮ ॥}
\cquote{ಹೇ ಪಾರ್ಥ, ಅಭ್ಯಾಸದ ಬಲದಿಂದ ಏಕಾಗ್ರವಾದ ಮನಸ್ಸಿನಿಂದ ಪರಮಪುರುಷನನ್ನು ಧ್ಯಾನ ಮಾಡಿದವನು ಈ ಭಗವಂತನನ್ನು ಹೊಂದುತ್ತಾನೆ.}
\slcol{\Index{ಕವಿಂ ಪುರಾಣಮನುಶಾಸಿತಾರಮ}ಣೋರಣೀಯಾoಸಮನುಸ್ಮರೇದ್ಯಃ ।\\
ಸರ್ವಸ್ಯ ಧಾತಾರಮಚಿಂತ್ಯರೂಪಮಾದಿತ್ಯವರ್ಣಂ ತಮಸಃ ಪರಸ್ತಾತ್ ॥ ೯ ॥}
\cquote{ಎಲ್ಲವನ್ನೂ ಬಲ್ಲ, ಸನಾತನವಾದ, ಜಗತ್ತನ್ನೆಲ್ಲ ನಿಯಂತ್ರಿಸುವ, ಸೂಕ್ಷ್ಮಕ್ಕಿಂತಲೂ ಸೂಕ್ಷ್ಮವಾದ, ಎಲ್ಲದರ ಧಾರಕನಾದ, ಅಚಿಂತ್ಯ ರೂಪನಾದ, ಸೂರ್ಯನಂತೆ ಬೆಳಗುವ, ಅಪ್ರಾಕೃತರೂಪನಾದ ಭಗವಂತನನ್ನು ಧ್ಯಾನಿಸಬೇಕು.}
\slcol{\Index{ಪ್ರಯಾಣಕಾಲೇ ಮನಸಾಚಲೇನ} ಭಕ್ತ್ಯಾ \\
\hspace*{0.5cm}ಯುಕ್ತೋ ಯೋಗಬಲೇನ ಚೈವ ।\\
ಭ್ರುವೋರ್ಮಧ್ಯೇ ಪ್ರಾಣಮಾವೇಶ್ಯ ಸಮ್ಯಕ್ಸ \\
\hspace*{0.5cm}ತಂ ಪರಂ ಪುರುಷಮುಪೈತಿ ದಿವ್ಯಮ್ ॥ ೧೦ ॥}
\cquote{ಸಾಯುವ ಕ್ಷಣದಲ್ಲಿ ಮನಸ್ಸನ್ನು ಸ್ಥಿರಗೊಳಿಸಿ, ಭಕ್ತಿಯುಕ್ತನಾಗಿ ಯೋಗ ಬಲದಿಂದ ಪ್ರಾಣ ವಾಯುವನ್ನು ಎರಡು ಹುಬ್ಬುಗಳ ನಡುವೆ ನೆಲೆಗೊಳಿಸಿದವನು ಆ ದಿವ್ಯವಾದ ಪರಮ ಪುರುಷನನ್ನು ಪಡೆಯುತ್ತಾನೆ.}
\slcol{\Index{ಯದಕ್ಷರಂ ವೇದವಿದೋ} ವದಂತಿ \\
\hspace*{0.5cm}ವಿಶಂತಿ ಯದ್ಯತಯೋ ವೀತರಾಗಾಃ ।\\
ಯದಿಚ್ಛಂತೋ ಬ್ರಹ್ಮಚರ್ಯಂ ಚರಂತಿ \\
\hspace*{0.5cm}ತತ್ತೇ ಪದಂ ಸಂಗ್ರಹೇಣ ಪ್ರವಕ್ಷ್ಯೇ ॥ ೧೧ ॥}
\cquote{ವೇದಗಳನ್ನರಿತವರು ಯಾವ ನಿತ್ಯಸ್ವರೂಪವನ್ನು ವರ್ಣಿಸುವರೋ, ಸಂಸಾರ ಶೋಕವನ್ನು ದಾಟಿದ ಯತಿಗಳು ಯಾವ ರೂಪವನ್ನು ಸೇರುವರೋ, ಯಾವುದನ್ನು ಬಯಸುವವರು ಮನಸ್ಸನ್ನು ಭಗವಂತನಲ್ಲಿ ನೆಲೆಗೊಳಿಸಿರುವರೋ ಅಂಥ ಭಗವದ್ರೂಪವನ್ನು ನಿನಗೆ ಸಂಕ್ಷೇಪವಾಗಿ ಹೇಳುತ್ತೇನೆ.}
\slcol{\Index{ಸರ್ವದ್ವಾರಾಣಿ ಸಂಯಮ್ಯ} ಮನೋ ಹೃದಿ ನಿರುಧ್ಯ ಚ ।\\
ಮೂರ್ಧ್ನ್ಯಾಧಾಯಾತ್ಮನಃ ಪ್ರಾಣಮಾಸ್ಥಿತೋ ಯೋಗಧಾರಣಾಮ್ ॥ ೧೨ ॥}
\cquote{ಮನಸ್ಸನ್ನು ಹೃದಯದಲ್ಲಿ ನಿಗ್ರಹಿಸಿ, ಎಲ್ಲಾ ಬಾಗಿಲುಗಳನ್ನು ನಿಗ್ರಹಿಸಿ, ಯೋಗ ಸ್ಥಿತಿಯಲ್ಲಿ ಪ್ರಾಣ ವಾಯುವನ್ನು ಎರಡು ಹುಬ್ಬುಗಳ ನಡುವೆ ನೆಲಗೊಳಿಸಿದವನು, ತನ್ನ ಆತ್ಮವನ್ನು ತಲೆಯ ಮೇಲೆ ನೆಲಗೊಳಿಸಿದವನು, ಆ ದಿವ್ಯವಾದ ಪರಮ ಪುರುಷನನ್ನು ಪಡೆಯುತ್ತಾನೆ.}
\slcol{\Index{ಓಮಿತ್ಯೇಕಾಕ್ಷರಂ ಬ್ರಹ್ಮ} ವ್ಯಾಹರನ್ಮಾಮನುಸ್ಮರನ್ ।\\
ಯಃ ಪ್ರಯಾತಿ ತ್ಯಜಂದೇಹಂ ಸ ಯಾತಿ ಪರಮಾಂ ಗತಿಮ್ ॥ ೧೩ ॥}
\cquote{ದೇಹದ ಎಲ್ಲಾ ಬಾಗಿಲುಗಳನ್ನು ನಿಗ್ರಹಿಸಿ, ಮನಸ್ಸನ್ನು ಭಗವಂತನಲ್ಲಿ ನೆಲಗೊಳಿಸಿ, ತನ್ನ ಪ್ರಾಣ ವಾಯುವನ್ನು ನಡುನೆತ್ತಿಗೇರಿಸಿ, ಯೋಗಸ್ಥಿತನಾಗಿ, ಬ್ರಹ್ಮದ ಹೆಸರಾದ “ಓಂ” ಎಂಬ ಒಂದಕ್ಷರವನ್ನು ಉಚ್ಚರಿಸುತ್ತ, ನನ್ನನ್ನು ನೆನೆಯುತ್ತ ಯಾವನು ದೇಹವನ್ನು ಬಿಟ್ಟು ಹೊರ ಬೀಳುತ್ತಾನೋ ಅವನು ಉತ್ತಮ ಗತಿಯನ್ನು ಹೊಂದುತ್ತಾನೆ.}
\slcol{\Index{ಅನನ್ಯಚೇತಾಃ ಸತತಂ} ಯೋ ಮಾಂ ಸ್ಮರತಿ ನಿತ್ಯಶಃ ।\\
ತಸ್ಯಾಹಂ ಸುಲಭಃ ಪಾರ್ಥ ನಿತ್ಯಯುಕ್ತಸ್ಯ ಯೋಗಿನಃ ॥ ೧೪ ॥}
\cquote{ಪಾರ್ಥ, ಯಾವನು ಬೇರೆ ಕಡೆಗೆ ಮನಸ್ಸನ್ನು ಹೋಗಗೋಡದೆ ನನ್ನನ್ನು ಎಡಬಿಡದೆ ಯಾವಾಗಲೂ ನೆನೆಯುತ್ತಾನೋ, ಯಾವಾಗಲೂ ಮನಸ್ಸನ್ನು ಬಿಗಿಹಿಡಿದಿರುವ ಯೋಗಿಯಾದ ಅವನಿಗೆ ನಾನು ಅನಾಯಾಸವಾಗಿ ಕೈ ಸೇರುವೆನು.}
\slcol{\Index{ಮಾಮುಪೇತ್ಯ ಪುನರ್ಜನ್ಮ} ದುಃಖಾಲಯಮಶಾಶ್ವತಮ್ ।\\
ನಾಪ್ನುವಂತಿ ಮಹಾತ್ಮಾನಃ ಸಂಸಿದ್ಧಿಂ ಪರಮಾಂ ಗತಾಃ ॥ ೧೫ ॥}
\cquote{ನನ್ನನ್ನು ಪಡೆದವರೆಂದರೆ ಸಾಧನೆಯ ಕೊನೆಯ ಸಿದ್ಧಿಯನ್ನು ಪಡೆದವರು. ಅಂಥವರು ದುಃಖದ ತವರುಮನೆಯಾದ ಮೂರು ದಿನದ ಬಾಳಿನ ಸಂಸಾರಕ್ಕೆ ಮತ್ತೆ ಮರಳಿ ಬರುವುದಿಲ್ಲ.}

\newpage
\begin{mananam}{\mananamTitle{ಮನನ ಶ್ಲೋಕ - ೧೦, ೧೩}}
\nfntmtxt ದೈಹಿಕ ಮತ್ತು ಮಾನಸಿಕ ಯೋಗಕ್ಷೇಮಕ್ಕಾಗಿ ವಿವಿಧ ವೈದ್ಯಕೀಯ ವಿಜ್ಞಾನಗಳಿರುವಂತೆಯೇ, ಭೌತಿಕ ದೇಹದಿಂದ ನಿರ್ಗಮಿಸುವ ಆತ್ಮಕ್ಕೆ ಮಾರ್ಗವನ್ನು ಬಹಿರಂಗಪಡಿಸುವ  ಆದರೆ, ಅಷ್ಟು ಪ್ರಸಿದ್ಧವಲ್ಲದ ವಿಜ್ಞಾನಗಳಿರಬಹುದೇ ? 
ನನ್ನ ಅನುಮಾನಗಳು ಮತ್ತು ಮರಣಾನಂತರದ ಜೀವನ ಹಾಗೂ ಪುನರ್ಜನ್ಮದ ಬಗ್ಗೆ ನಂಬಿಕೆಯ ಕೊರತೆಯೇ ನಾನು ಸ್ವತಃ ಯಾವುದೇ ವಿಚಾರಣೆಯನ್ನು ಕೈಗೊಳ್ಳದಂತೆ ತಡೆಯುತ್ತವೆಯೇ?  ಈ ದೇಹ-ಮನಸ್ಸನ್ನು ಮೀರಿ ಆತ್ಮವಿದೆಯೇ ಎಂಬ ಬಗ್ಗೆ ವೈಯುಕ್ತಿಕ ಅನುಭವ ಅಥವಾ ತಿಳುವಳಿಕೆಯನ್ನು ಪಡೆಯಲು ನಾನು ಯಾವುದೇ ಗಂಭೀರ ವಿಚಾರಣೆಯ ಅಭ್ಯಾಸ ಮಾಡಿದ್ದೇನೆಯೇ ಅಥವಾ ವಿಚಾರಣೆಯನ್ನು ಕೈಗೊಂಡಿದ್ದೇನೆಯೇ?  ನಾವು ಪರಿಕಲ್ಪನಾತ್ಮಕವಾಗಿ “ದೇವರು” ಎಂದು ಕರೆಯುವ, ಒಂದು ಸರ್ವೋಚ್ಚ ಅಸ್ತಿತ್ವ ಹಾಗೂ ಅದರ ಸ್ವರೂಪವಿದೆಯೇ ಎಂದು ನಾನು ಪರಿಶೋಧಿಸಿದ್ದೇನೆಯೇ?
\end{mananam}
\sReflect
\begin{inspiration}{\mananamTitle{ಸ್ಪೂರ್ತಿ}}
\nfntmtxt \footnotesize ಭಗವಂತ ತನ್ನನ್ನು ಒಲಿಸಿಕೊಳ್ಳಲು ಒಂದು ನಿರ್ದಿಷ್ಟವಾದ ತಂತ್ರವನ್ನು ಒದಗಿಸುತ್ತಾನೆ ಏಕೆಂದರೆ, ಹೆಚ್ಚಿನವರಿಗೆ ಕೇವಲ ಪದಗಳು ಸಾಕಾಗುವುದಿಲ್ಲ. ಈ ತಂತ್ರವು, “ಓಂ” ಎಂಬ ಪವಿತ್ರವಾದ ಶಬ್ದವನ್ನು ಉಪಯೋಗಿಸಿಕೊಂಡು, ಮನಸ್ಸನ್ನು ನಿಶ್ಚಲಗೊಳಿಸಿ, ದೇಹದಲ್ಲಿ ಎರಡು ಹುಬ್ಬುಗಳ ನಡುವೆ ಇರುವ ಆಧ್ಯಾತ್ಮಿಕ ಶಕ್ತಿ ಕೇಂದ್ರದ ಮೇಲೆ ಗಮನವನ್ನು ಕೇಂದ್ರೀಕರಿಸುವುದನ್ನು ಒಳಗೊಂಡಿರುತ್ತದೆ. ಸಾವಿನ ಕ್ಷಣದಲ್ಲಿ ಈ ತಂತ್ರವನ್ನು ಪರಿಣಾಮಕಾರಿಯಾಗಿ ಉಪಯೋಗಿಸಿಕೊಳ್ಳಲು ಜೀವನವಿಡೀ ಅದನ್ನು ಅಭ್ಯಾಸ ಮಾಡುತ್ತಿರಬೇಕು ಮತ್ತು ಅದರ ಪರಿಣಾಮವನ್ನು ಪರಿಶೀಲಿಸುತ್ತಿರಬೇಕು.
ಒಬ್ಬ ಯೋಧನು ನಿಯಮಿತವಾಗಿ ಅಭ್ಯಾಸ ಮಾಡಿ ಯುದ್ಧಕ್ಕೆ ಸನ್ನದ್ಧನಾಗುವಂತೆಯೇ, ನಾವು ಜೀವನದಲ್ಲಿ ಸತತ ಅಭ್ಯಾಸದ ಮೂಲಕ ನಮ್ಮನ್ನು ನಾವು ಸಿದ್ದಪಡಿಸಿಕೊಳ್ಳುವುದರಿಂದ ನಾವು ಸಾವನ್ನು ಎದುರಿಸಲು ಸಿದ್ದರಾಗುತ್ತೇವೆ. 
\end{inspiration}
\newpage

\slcol{\Index{ಆಬ್ರಹ್ಮಭುವನಾಲ್ಲೋಕಾಃ} ಪುನರಾವರ್ತಿನೋಽರ್ಜುನ ।\\
ಮಾಮುಪೇತ್ಯ ತು ಕೌಂತೇಯ ಪುನರ್ಜನ್ಮ ನ ವಿದ್ಯತೇ ॥ ೧೬ ॥}
\cquote{ಅರ್ಜುನ, ಬ್ರಹ್ಮ ಲೋಕವನ್ನು ಮೊದಲು ಮಾಡಿಕೊಂಡು ಎಲ್ಲಾ ಲೋಕಗಳನ್ನು ಪಡೆದಿರುವವರು ಮತ್ತೆ ಹುಟ್ಟುವವರೇ. ಕೌಂತೇಯ, ನನ್ನನ್ನು ಪಡೆದರೆ ಮತ್ತೆ ಹುಟ್ಟುವಿಕೆ ಇಲ್ಲ.}
\slcol{\Index{ಸಹಸ್ರಯುಗಪರ್ಯಂತಮ}ಹರ್ಯದ್ಬ್ರಹ್ಮಣೋ ವಿದುಃ ।\\
ರಾತ್ರಿಂ ಯುಗಸಹಸ್ರಾಂತಾಂ ತೇಽಹೋರಾತ್ರವಿದೋ ಜನಾಃ ॥ ೧೭ ॥}
\cquote{ಕಾಲ ಪರಿಮಾಣವನ್ನರಿತ ಜನರು ಸಾವಿರಾರು ಯುಗಗಳು ಭಗವಂತನ ಒಂದು ಹಗಲೆಂದೂ ಹಾಗೆಯೇ ಸಾವಿರಾರು ಯುಗಗಳನ್ನು ಇರುಳೆಂದೂ ತಿಳಿಯುವವರು.}
\slcol{\Index{ಅವ್ಯಕ್ತಾದ್ವ್ಯಕ್ತಯಃ ಸರ್ವಾಃ} ಪ್ರಭವಂತ್ಯಹರಾಗಮೇ ।\\
ರಾತ್ರ್ಯಾಗಮೇ ಪ್ರಲೀಯಂತೇ ತತ್ರೈವಾವ್ಯಕ್ತಸಂಙ್ಞಕೇ ॥ ೧೮ ॥}
\cquote{ಹಗಲು ಬರುವಾಗ ಚರಾಚರಗಳೆಲ್ಲ ಅವ್ಯಕ್ತನಾದ ಭಗವಂತನಿಂದ ಹುಟ್ಟುತ್ತವೆ. ಇರುಳಾದಾಗ ಅವು ಆ ಅವ್ಯಕ್ತದಲ್ಲಿಯೇ ಅಡಗುತ್ತವೆ.}
\slcol{\Index{ಭೂತಗ್ರಾಮಃ ಸ ಏವಾಯಂ} ಭೂತ್ವಾ ಭೂತ್ವಾ ಪ್ರಲೀಯತೇ ।\\
ರಾತ್ರ್ಯಾಗಮೇಽವಶಃ ಪಾರ್ಥ ಪ್ರಭವತ್ಯಹರಾಗಮೇ ॥ ೧೯ ॥}
\cquote{ಪಾರ್ಥ, ಹಿಂದಿದ್ದ ಅದೇ ಈ ಚರಾಚರವು ಹುಟ್ಟಿ ಹುಟ್ಟಿ ಇರುಳಾಗುತ್ತಲೇ ತನಗರಿವಿಲ್ಲದೆ ಅಡಗುತ್ತದೆ, ಹಗಲು ಬರುತ್ತಲೇ ಹುಟ್ಟುತ್ತದೆ.}
\slcol{\Index{ಪರಸ್ತಸ್ಮಾತ್ತು ಭಾವೋ}ಽನ್ಯೋಽವ್ಯಕ್ತೋಽವ್ಯಕ್ತಾತ್ಸನಾತನಃ ।\\
ಯಃ ಸ ಸರ್ವೇಷು ಭೂತೇಷು ನಶ್ಯತ್ಸು ನ ವಿನಶ್ಯತಿ ॥ ೨೦ ॥}
\cquote{ಯಾವನು ಆ ಜಗತ್ತಿನ ಬೀಜಕ್ಕಿಂತಲೂ ಆಚೆಗೆ ಅಗೋಚರನಾಗಿರುವನೋ ಅವನು ಎಲ್ಲ ಚರಾಚರ  ವಸ್ತುಗಳ ನಾಶವಾದರೂ ನಾಶವಾಗುವುದಿಲ್ಲ.}

\newpage
\slcol{\Index{ಅವ್ಯಕ್ತೋಽಕ್ಷರ ಇತ್ಯುಕ್ತಸ್ತಮಾಹುಃ} ಪರಮಾಂ ಗತಿಮ್ ।\\
ಯಂ ಪ್ರಾಪ್ಯ ನ ನಿವರ್ತಂತೇ ತದ್ಧಾಮ ಪರಮಂ ಮಮ ॥ ೨೧ ॥}
\cquote{ಅವ್ಯಕ್ತನಾದ ಭಗವಂತನೇ ನಾಶರಹಿತನಾದ ಅಕ್ಷರ ಪುರುಷನು. ಅವನೇ “ಪರಮ ಗತಿ” ಎಂದು ಹೇಳುತ್ತಾರೆ. ಯಾವುದನ್ನು ಹೊಂದಿದ ಮೇಲೆ ಮತ್ತೆ ಹುಟ್ಟು ಸಾವುಗಳಿಗೆ ತಿರುಗುವುದಿಲ್ಲವೋ ಅದೇ ನನ್ನ ದಿವ್ಯವಾದ ಸ್ವರೂಪವು.}
\slcol{\Index{ಪುರುಷಃ ಸ ಪರಃ ಪಾರ್ಥ} ಭಕ್ತ್ಯಾ ಲಭ್ಯಸ್ತ್ವನನ್ಯಯಾ ।\\
ಯಸ್ಯಾಂತಃಸ್ಥಾನಿ ಭೂತಾನಿ ಯೇನ ಸರ್ವಮಿದಂ ತತಮ್ ॥ ೨೨ ॥}
\cquote{ಪಾರ್ಥ, ಚರಾಚರ ವಸ್ತುಗಳೆಲ್ಲ ಯಾವನ ಒಳಗಿರುವುವೋ ಇವೆಲ್ಲವೂ ಯಾರಿಂದ ತುಂಬಿಕೊಂಡಿವೆಯೋ ಆ ಪರಮ ಪುರುಷನನ್ನು ಅನನ್ಯ ಭಕ್ತಿಯಿಂದ ಮಾತ್ರವೇ ಪಡೆಯುವುದು ಸಾಧ್ಯ.}
\slcol{\Index{ಯತ್ರ ಕಾಲೇ ತ್ವನಾವೃತ್ತಿ}ಮಾವೃತ್ತಿಂ ಚೈವ ಯೋಗಿನಃ ।\\
ಪ್ರಯಾತಾ ಯಾಂತಿ ತಂ ಕಾಲಂ ವಕ್ಷ್ಯಾಮಿ ಭರತರ್ಷಭ ॥ ೨೩ ॥}
\cquote{ಭಾರತವೀರಾ, ದೇಹವನ್ನು ಬಿಟ್ಟು ಹೊರಡುವವರು ಯಾವ ಕಾಲಾಭಿಮಾನಿ ದೇವತೆಗಳ ದಾರಿಯಲ್ಲಿ ಹೋಗಿ ತಿರುಗಿ ಬರದೆ ಇರುವಿಕೆಯನ್ನೂ, ತಿರುಗಿ ಬರುವಿಕೆಯನ್ನೊ ಪಡೆಯುತ್ತಾರೋ, ಕಾಲಾಭಿಮಾನಿ ದೇವತೆಗಳಿರುವ ಆ ದಾರಿಯನ್ನು ಹೇಳುತ್ತೇನೆ.}
\slcol{\Index{ಅಗ್ನಿರ್ಜೋತಿರಹಃ ಶುಕ್ಲಃ} ಷಣ್ಮಾಸಾ ಉತ್ತರಾಯಣಮ್ ।\\
ತತ್ರ ಪ್ರಯಾತಾ ಗಚ್ಛಂತಿ ಬ್ರಹ್ಮ ಬ್ರಹ್ಮವಿದೋ ಜನಾಃ ॥ ೨೪ ॥}
\cquote{ಅಗ್ನಿ ದೇವತೆ, ಜ್ಯೋತಿಯ ದೇವತೆ, ಹಗಲಿನ ದೇವತೆ, ಶುಕ್ಲ ಪಕ್ಷದ ದೇವತೆ, ಉತ್ತರಾಯಣದ ಆರು ತಿಂಗಳ ದೇವತೆ- ಇವರಿರುವ ದಾರಿಯಲ್ಲಿ ಹೋದ ಬ್ರಹ್ಮ ಜ್ಞಾನಿಗಳು ಭಗವಂತನನ್ನು ಪಡೆಯುವರು.}
\slcol{\Index{ಧೂಮೋ ರಾತ್ರಿಸ್ತಥಾ ಕೃಷ್ಣಃ} ಷಣ್ಮಾಸಾ ದಕ್ಷಿಣಾಯನಮ್ ।\\
ತತ್ರ ಚಾಂದ್ರಮಸಂ ಜ್ಯೋತಿರ್ಯೋಗೀ ಪ್ರಾಪ್ಯ ನಿವರ್ತತೇ ॥ ೨೫ ॥}
\cquote{ಧೂಮದ ದೇವತೆ, ಇರುಳಿನ ದೇವತೆ, ಹಾಗೆಯೇ ಕೃಷ್ಣ ಪಕ್ಷದ ದೇವತೆ, ದಕ್ಷಿಣಾಯನದ ಆರು ತಿಂಗಳ ದೇವತೆ- ಇವರಿರುವ ದಾರಿಯಲ್ಲಿ ಹೋದ ಯೋಗಿಗಳು ಚಂದ್ರನ ಲೋಕವನ್ನು ಹೊಂದಿ ಮತ್ತೆ ಹಿಂದಿರುಗುತ್ತಾರೆ.}

\newpage
\begin{mananam}{\mananamTitle{ಮನನ ಶ್ಲೋಕ - 20, 21}}
\nfntmtxt ಎಲ್ಲಾ ಸೃಷ್ಟಿಯ ಬೀಜವು ಅಡಗಿರುವ ‘ಅವ್ಯಕ್ತ’ದ ಬಗ್ಗೆ ನನ್ನ ತಿಳುವಳಿಕೆ ಏನು? ಈ ಜಗತ್ತು ಎಂದಾದರೂ ಕೊನೆಗೊಳ್ಳುವ ಒಂದು ಸಮಯವಿರಬಹುದೇ? ಪ್ರತಿ ರಾತ್ರಿಯೂ ನಾನು ಸುಪ್ತಾವಸ್ಥೆಯಲ್ಲಿದ್ದಾಗ (ಆಳವಾದ ನಿದ್ರೆ), ಈ ಪ್ರಪಂಚದ ಬಗೆಗಿರುವ ನನ್ನ ಗ್ರಹಿಕೆಯು ವಿಸರ್ಜನೆಯಾಗುವುದನ್ನು ನಾನು ಅನುಭವಿಸುತ್ತೇನಲ್ಲವೇ? ಬೆಳಿಗ್ಗೆ ಎಚ್ಚರವಾದಾಗ, ನನಗೆ ಈ ಪ್ರಪಂಚವು ಪುನಃ ಹೇಗೆ ಹೊರಹೊಮ್ಮುತ್ತದೆ?   ನೆನಪುಗಳನ್ನು ಮೆಲುಕು ಹಾಕಲು, ಸ್ವಯಂ ಅರಿವನ್ನು ಮರಳಿ ಪಡೆಯಲು, ಮತ್ತು, ಮತ್ತೆ ಪ್ರಪಂಚದ ಅರಿವನ್ನು ಹೊಂದಲು ನನಗೆ ಹೇಗೆ ಸಾಧ್ಯವಾಗುತ್ತದೆ? ಗಾಢ ನಿದ್ರೆಯಲ್ಲಿದ್ದಾಗ “ನಾನು” (ಇರುವಿಕೆ) ಇದ್ದೆನೋ ಇಲ್ಲವೋ? ಅಂತೆಯೇ, ಹಿಂದೂ ಧರ್ಮಗ್ರಂಥಗಳು ಸೂಚಿಸುವಂತೆ, ಪ್ರಪಂಚದ ಸಾಮಾನ್ಯ ಅನುಭವವೂ, ವಿಸರ್ಜನೆ ಮತ್ತು ಪುನರುತ್ಥಾನಕ್ಕೆ ಒಳಗಾಗಬಹುದೇ?
\end{mananam}
\sReflect
\begin{inspiration}{\mananamTitle{ಸ್ಪೂರ್ತಿ}}
\nfntmtxt \footnotesize ಒಂದು ಇಡೀ ವೃಕ್ಷವು ಒಂದು ಸಣ್ಣ ಬೀಜದಿಂದ ಹೇಗೆ ಹೊರಹೊಮ್ಮುತ್ತದೆ ಎಂಬುದು ಪ್ರಪಂಚದ ಅತ್ಯಂತ ಅದ್ಭುತ ಮತ್ತು ರಹಸ್ಯಗಳಲ್ಲಿ ಒಂದಾಗಿದೆ. ಮರದ ಪೂರ್ಣ ವಿನ್ಯಾಸವು ಬೀಜದ ಒಳಗೆ ಅಡಕವಾಗಿರುವಂತೆಯೇ, ಇಡೀ ಬ್ರಹ್ಮಾಂಡದ ಮುದ್ರೆಯು ಸೃಷ್ಟಿಯಾಗುವ ಮೊದಲು ಅವ್ಯಕ್ತ ಸ್ಥಿತಿಯಲ್ಲಿರುತ್ತದೆ;ಪುನಃ ಲಯವಾಗುವಾಗ ಅವ್ಯಕ್ತಾವಸ್ಥೆಗೆ ತೆರಳುತ್ತದೆ. 
ಅವ್ಯಕ್ತವನ್ನು ಮೀರಿ ಉಳಿದ ಶಾಶ್ವತ ಪ್ರಜ್ಞೆ ಅಥವಾ ಅಸ್ತಿತ್ವವು ಎಲ್ಲಾ ಆತ್ಮಗಳ ಅಂತಿಮ ಗಮ್ಯವಾಗಿದೆ. ಆ ವಾಸ್ತವಕ್ಕೆ ನಮ್ಮ ಮಾರ್ಗವನ್ನು ಪುನಃ ಕಂಡುಕೊಳ್ಳುವುದು ನಮ್ಮನ್ನು ಅಂತಿಮ ವಿಶ್ರಾಂತಿ ತಾಣಕ್ಕೆ ಕರೆದೊಯ್ಯುತ್ತದೆ. ಯೋಗದ ಸಂಪೂರ್ಣ ವಿಜ್ಞಾನವು ಈ ಅತ್ಯಗತ್ಯ ವಾಸ್ತವಕ್ಕೆ ನಮಗೆ ಮರಳಿ ಮಾರ್ಗದರ್ಶನ ನೀಡುವ ಗುರಿಯನ್ನು ಹೊಂದಿದ್ದು, ಜನನ-ಮರಣದ ಚಕ್ರಗಳಿಂದ  ನಮ್ಮನ್ನು ಮುಕ್ತಗೊಳಿಸುತ್ತದೆ.
\end{inspiration}
\newpage

\begin{mananam}{\mananamTitle{ಮನನ ಶ್ಲೋಕ - 22}}
\nfntmtxt ಅವ್ಯಕ್ತ ಸ್ಥಿತಿ ಮತ್ತು ಅದನ್ನು ಮೀರಿದ ಪರಮಾತ್ಮನನ್ನು ಕುರಿತು ನಾನೇಕೆ ವಿಚಾರಣೆಯಲ್ಲಿ ತೊಡಗಬೇಕು? “ನೋಡುವುದನ್ನೇ ನಂಬುವುದು” ಎಂಬ  ಸಿದ್ಧಾಂತವನ್ನು ಅನುಸರಿಸುವ ಬಹುಸಂಖ್ಯಾತರಂತೆ ನಾನೂ ಸಹ ತೃಪ್ತಿ ಹೊಂದಲು ಸಾಧ್ಯವಿಲ್ಲವೇ? ಅವ್ಯಕ್ತ ಸ್ಥಿತಿ ಮತ್ತು ಪರಮಾತ್ಮನನ್ನು ಅರಿಯುವುದರಿಂದ ನನ್ನ ಬಗ್ಗೆಯೇ ಇನ್ನಷ್ಟು ತಿಳಿದುಕೊಳ್ಳಲು ಹೇಗೆ ಸಹಾಯವಾಗುತ್ತದೆ? ಸಕಲ ಸೃಷ್ಟಿಯ ಅತ್ಯಗತ್ಯ ಮೂಲವು ಈ ಸೃಷ್ಟಿಯಲ್ಲಿಯೇ ಇರಬಹುದೇ? ನೋಡುವ, ಕೇಳುವ, ಯೋಚಿಸುವ, ಉಸಿರಾಡುವ,ಹಾಗೂ ಜೀರ್ಣ ಮಾಡಿಕೊಳ್ಳುವ ನನ್ನ ಸಾಮರ್ಥ್ಯಗಳ ಹಿಂದಿನ ನಿಜವಾದ ಶಕ್ತಿ ಏನು? ಈ ಭೌತಿಕ ದೇಹಕ್ಕೆ ಉಸಿರಿನ ಅಗತ್ಯವಿರುವಂತೆಯೇ, ಎಲ್ಲಾ ಜೀವಿಗಳಿಗೂ ಸಾಮಾನ್ಯವಾದ ಒಂದು ಅಗತ್ಯ ಗುಣಲಕ್ಷಣವಿದೆಯೇ? 
\end{mananam}
\sReflect
\begin{inspiration}{\mananamTitle{ಸ್ಪೂರ್ತಿ}}
\nfntmtxt ನಾವು ನಮ್ಮ ಕುಟುಂಬದವರಿಗೆ ಮತ್ತು ಸ್ನೇಹಿತರಿಗೆ ನಮ್ಮ ಪ್ರೀತಿ ಮತ್ತು ಸಮರ್ಪಣೆಯನ್ನು ಎಷ್ಟು ಸುಲಭವಾಗಿ ನೀಡುತ್ತೇವೆ; ನಾವು ಎಷ್ಟು ಸ್ವಾಭಾವಿಕವಾಗಿ ನಮ್ಮ  ಮೇಲಾಧಿಕಾರಿಗಳು ಮತ್ತು ಅಧಿಕಾರದಲ್ಲಿರುವವರ ಕಡೆಗೆ ನಮ್ಮ ನಿಷ್ಠೆಯನ್ನು ತಿರುಗಿಸುತ್ತೇವೆ; ಆದರೂ, ಎಲ್ಲಾ ಸೃಷ್ಟಿಯ ಅತ್ಯಗತ್ಯ ಮೂಲದ ಮೇಲೆ ನಾವು ಎಷ್ಟು ಅಲ್ಪವಾಗಿ ಗಮನ ಹರಿಸುತ್ತೇವೆ!
ನಮ್ಮದೇ ಆದಂತಹ ಜಾಗೃತಿ - ಅದರ ಮೂಲಕ ನಾವು ಈ ಜಗತ್ತನ್ನು ಅನುಭವಿಸುತ್ತಾ, ಅದರಲ್ಲಿ ಸಾಗುತ್ತಾ, ನಮ್ಮ ಆಲೋಚನೆಗಳನ್ನು ಯೋಚಿಸುತ್ತೇವೋ — ಅದು ಶಕ್ತಿ ಮತ್ತು ಪ್ರಜ್ಞೆಯಿಂದ ತುಂಬಲ್ಪಟ್ಟಿದೆ; ಇದೇ ಸಾಮಾನ್ಯ ಮಾನವನ ಅಜ್ಞಾನ— ಇದರಿಂದ ಹೊರಬರಲು ಧರ್ಮಗ್ರಂಥಗಳು ಮತ್ತು ಋಷಿಮುನಿಗಳು ನಮ್ಮನ್ನು ಪ್ರಚೋದಿಸುತ್ತಾರೆ; ನನ್ನ ಕಣ್ಣು ಮತ್ತು ಇಂದ್ರಿಯಗಳಿಗೆ ಗೋಚರಿಸದಿದ್ದರೂ, ಇಡೀ ಸೃಷ್ಟಿಯನ್ನು ವ್ಯಾಪಿಸಿರುವ ಆ ತತ್ವವನ್ನು ಅರ್ಥಮಾಡಿಕೊಳ್ಳಲು ಅವರು ನಮ್ಮನ್ನು ತೀವ್ರವಾಗಿ ಪ್ರೆರೇಪಿಸುತ್ತಾರೆ. 
\end{inspiration}
\newpage

\slcol{\Index{ಶುಕ್ಲಕೃಷ್ಣೇ ಗತೀ ಹ್ಯೇತೇ} ಜಗತಃ ಶಾಶ್ವತೇ ಮತೇ ।\\
ಏಕಯಾ ಯಾತ್ಯನಾವೃತ್ತಿಮನ್ಯಯಾವರ್ತತೇ ಪುನಃ ॥ ೨೬ ॥}
\cquote{ಈ ಬೆಳಕಿನ ಮತ್ತು ಕತ್ತಲೆಯ ದಾರಿಗಳೆರಡು ಜಗತ್ತಿನಲ್ಲಿ ಯಾವಾಗಲೂ ಇರುವವುಗಳೆನಿಸುವುವು. ಒಂದರಿಂದ ಹೋದವನು ತಿರುಗಿ ಬರುವುದಿಲ್ಲ, ಇನ್ನೊಂದರಿಂದ ಹೋದವನು ಮತ್ತೆ ತಿರುಗಿ ಬರುತ್ತಾನೆ.}
\slcol{\Index{ನೈತೇ ಸೃತೀ ಪಾರ್ಥ} ಜಾನನ್ಯೋಗೀ ಮುಹ್ಯತಿ ಕಶ್ಚನ ।\\
ತಸ್ಮಾತ್ಸರ್ವೇಷು ಕಾಲೇಷು ಯೋಗಯುಕ್ತೋ ಭವಾರ್ಜುನ ॥ ೨೭ ॥}
\cquote{ಪಾರ್ಥ, ಈ ಎರಡೂ ದಾರಿಗಳನ್ನು ತಿಳಿದ ಯೋಗಿಯು ಮೋಹಗೊಳ್ಳುವುದಿಲ್ಲ. ಅರ್ಜುನ, ಆದ್ದರಿಂದ ಎಲ್ಲ ಕಾಲದಲ್ಲಿಯೂ ಧ್ಯಾನ ಯೋಗವನ್ನು ನಡೆಸುತ್ತಿರು.}
\slcol{\Index{ವೇದೇಷು ಯಜ್ಞೇಷು} ತಪಃಸು ಚೈವ \\
\hspace*{0.5cm}ದಾನೇಷು ಯತ್ಪುಣ್ಯಫಲಂ ಪ್ರದಿಷ್ಟಮ್ ।\\
ಅತ್ಯೇತಿ ತತ್ಸರ್ವಮಿದಂ ವಿದಿತ್ವಾಯೋಗೀ \\
\hspace*{0.5cm}ಪರಂ ಸ್ಥಾನಮುಪೈತಿ ಚಾದ್ಯಮ್ ॥ ೨೯ ॥}
\cquote{ವೇದಗಳನ್ನು ಓದಿದರೆ, ಯಜ್ಞಗಳನ್ನು ಮಾಡಿದರೆ, ಉಪವಾಸ ಮೊದಲಾದ ತಪಸ್ಸನ್ನು ನಡೆಸಿದರೆ, ದಾನ ಮಾಡಿದರೆ ಯಾವ ಫಲವು ಶಾಸ್ತ್ರದಲ್ಲಿ ಹೇಳಿದೆಯೋ ಅದೆಲ್ಲಕ್ಕಿಂತ ಮಿಗಿಲಾದ ಫಲವನ್ನು ಇದನ್ನರಿತ ಜ್ಞಾನ ಯೋಗಿ ಪಡೆಯುತ್ತಾನೆ ಮತ್ತು ಎಲ್ಲಕ್ಕಿಂತ ಶ್ರೇಷ್ಠವಾದ ಪರಮ ಪದವನ್ನು ಪಡೆಯುತ್ತಾನೆ.}

\newpage
\begin{mananam}{\mananamTitle{ಮನನ ಶ್ಲೋಕ - ೨೮}}
\nfntmtxt \footnotesize ಮರಣಾನಂತರದ ಜೀವನದ ಕಲ್ಪನೆಗಳ ಆಧಾರದ ಮೇಲೆ, ಈ ಜೀವನದಲ್ಲಿ ನನಗೆ ಲಭ್ಯವಿರುವ ಜೀವನದ ಆಯ್ಕೆಗಳನ್ನು ನಾನು ಹೇಗೆ ಬಳಸುತ್ತಿದ್ದೇನೆ? ನನ್ನ ಮರಣಾನಂತರದ ಜೀವನಕ್ಕೆ ಮತ್ತು ನನ್ನ ಪೂರ್ವಜರಿಗೆ ಪ್ರಯೋಜನಕಾರಿ ಎಂದು ಹೇಳಲಾಗುವ ವೈದಿಕ ಆಜ್ಞೆಗಳಿಗೆ ಸಂಬಂಧಿಸಿದ ವಿಧಿಗಳು, ಆಚರಣೆಗಳು ಮತ್ತು ಯಜ್ಞಗಳನ್ನು ನಾನು ಎಷ್ಟರ ಮಟ್ಟಿಗೆ ನಂಬುತ್ತೇನೆ ಮತ್ತು ಅಭ್ಯಾಸ ಮಾಡುತ್ತೇನೆ? 
ನಾನು ಕೇವಲ ನನ್ನ ಯೋಗಕ್ಷೇಮಕ್ಕಾಗಿ  ಪರೋಪಕಾರ ಮತ್ತು ಒಳ್ಳೆಯ ಕಾರ್ಯಗಳಲ್ಲಿ ಭಾಗವಹಿಸುತ್ತೇನೆಯೋ ಅಥವಾ ಅದರಿಂದ ಲಭಿಸುವ ಪುಣ್ಯವು ನನಗೆ ಜನ್ಮಾಂತರಕ್ಕೆ ಒದಗಿ ಬರುವುದೆಂಬ ಭರವಸೆಯಿಂದಲೋ? ಮರಣದ ಘಳಿಗೆಯಲ್ಲಿ ನನಗೆ ಸಹಾಯ ಮಾಡುವಂತಹ ಯಾವ ಜ್ಞಾನ ಅಥವಾ ತಿಳುವಳಿಕೆಯನ್ನು ನಾನು ಪಡೆದುಕೊಂಡಿದ್ದೇನೆ, ಅಥವಾ ಅದಕ್ಕಾಗಿ, ಈಗ ಯಾವ ಅಭ್ಯಾಸಗಳನ್ನು ಮಾಡುತ್ತಿದ್ದೇನೆ? ಸಾವಿನ ಕ್ಷಣ ಹಾಗೂ ಸಾವಿನಾಚೆಗೂ ತಯಾರಾಗಲು, ನನಗೆ ಸಹಾಯ ಮಾಡುವಂತಹ ಯಾವುದೇ ಧರ್ಮಗ್ರಂಥ ಆಧಾರಿತ ವಿವೇಕ ಅಥವಾ ಅಭ್ಯಾಸಗಳು ಇವೆಯೇ?
\end{mananam}
\sReflect
\begin{inspiration}{\mananamTitle{ಸ್ಪೂರ್ತಿ}}
\nfntmtxt \footnotesize ಈ ಜೀವನದಲ್ಲಿ ಸದಾ ನಮಗೆ ಎರಡು ಆಯ್ಕೆಯ ಮಾರ್ಗಗಳು ಹೇಗೆ ಪ್ರಸ್ತುತಪಡಿಸಲ್ಪಡುತ್ತವೆಯೋ  ಹಾಗೆಯೇ, ಸಾವಿನ ಕ್ಷಣದಲ್ಲಿ ನಿರ್ಗಮಿಸುವ ಪ್ರತಿಯೊಂದು ಆತ್ಮವೂ ಆಯ್ಕೆ ಮಾಡಲು ಎರಡು ಮಾರ್ಗಗಳನ್ನು ಹೊಂದಿದೆ: ಒಂದು  ಮಾರ್ಗ ಉನ್ನತ ಅಸ್ತಿತ್ವದೆಡೆಗೆ ಕೊಂಡೊಯ್ಯುತ್ತದೆಯಾದರೆ  ಮತ್ತೊಂದು,  ಅವರವರ ಪುಣ್ಯ ಮತ್ತು ಪಾಪದ ಸಂಗ್ರಹಕ್ಕೆ ಅನುಸಾರವಾಗಿ ಪುನಃ ಭೂಮಿಯ ಮೇಲೆ ಅಥವಾ ಇನ್ನೂ ಕೆಳ ಅಸ್ತಿತ್ವದ ಲೋಕದಲ್ಲಿ ಜನ್ಮ ತಾಳುವುದಾಗುತ್ತದೆ. ಪುಣ್ಯದ ಮಾರ್ಗವನ್ನು ಅನುಸರಿಸಲು ಸಹಾಯ ಮಾಡಲು, ವೇದಗಳು ವಿವಿಧ ವೈದೀಕ ಆಚರಣೆಗಳು ಮತ್ತು ಯಜ್ಞಗಳನ್ನು ಒದಗಿಸುತ್ತವೆ. ಆದಾಗ್ಯೂ, ಯೋಗಿಯ ಗುರಿ ಈ ಎರಡೂ ಮಾರ್ಗಗಳಲ್ಲ, ಏಕೆಂದರೆ ಯಾವುದೇ ಆಯ್ಕೆಯೂ ಇನ್ನೂ ಮಿತಿಗೆ ಬದ್ಧವಾಗಿರುತ್ತದೆ. 
ಸತ್ಯದ ಸಾಕ್ಷಾತ್ಕಾರದಿಂದ ಮತ್ತು ಸತ್ಯದಲ್ಲಿ ಬದ್ಧನಾಗಿರುವ ಮೂಲಕ ಒಬ್ಬ ಯೋಗಿಯು, ಆ ‘ಅಕ್ಷರ ಬ್ರಹ್ಮ’ನ ಸರ್ವೋಚ್ಚ ಸ್ಥಿತಿಯನ್ನು ಪಡೆದು, ಸೃಷ್ಟಿಯಲ್ಲವನ್ನೂ ಮೀರಿದ, ಪಾಪ-ಪುಣ್ಯಗಳ ಸಂಗ್ರಹವನ್ನೂ ಮೀರಿದ, ‘ಸಂಪೂರ್ಣತ್ವವನ್ನು’ ಕಂಡುಕೊಂಡ ನಂತರ ಬೇರೇನನ್ನೂ ಹುಡುಕುವುದಿಲ್ಲ.
\end{inspiration}
\newpage

\chapEndSloka{ಅಕ್ಷರ ಬ್ರಹ್ಮ ಯೋಗ}